%=============================================================================
\section{Methods}
\label{sec:methods}
%=============================================================================

\subsection{Dataset provenance}
\label{subsec:dataset}

We interrogated the American Gut Project (\agp)~\citep{mcdonald2018american},
a citizen-science 16S rRNA sequencing initiative whose public BIOM release
embodies one of the largest annotated gut microbiome collections available
without access controls.
Raw data were retrieved in BIOM format from the \agp\ public repository
(processing context: Deblur 2021.09, Illumina 16S V4, 150\,nt amplicons,
\texttt{Deblur\_2021.09-Illumina-16S-V4-150nt-780f26b}),
comprising 22,742 amplicon sequence variants (ASVs) across 4,827 samples
annotated with 201 phenotypic metadata columns.

All analyses were restricted to fecal samples (\texttt{BODY\_SITE =
UBERON:feces}; $n = 3{,}679$ pre-filtration), the single body site with
sufficient representation for stratified group comparisons.
For the \ibd\ comparison, the healthy group comprises samples whose
respondents explicitly selected ``I do not have IBD'' ($n = 3{,}084$
post-filtration); approximately 160 fecal samples were excluded from
that comparison because their IBD field was blank or ambiguous.
These samples are retained for the diet and antibiotic comparisons,
where IBD status is not a filter criterion.

\subsection{Preprocessing}
\label{subsec:preprocessing}

Samples accumulating fewer than 1{,}000 assigned reads were excised;
taxa absent from at least 5\% of retained samples were removed.
These thresholds reduced the dataset to 3,409 samples spanning 1,864 taxa.

Compositional distortion inherent in proportional sequencing data was
corrected via the centered log-ratio (\clr) transformation~\citep{aitchison1982statistical,
gloor2017microbiome}:
\begin{equation}
  \operatorname{CLR}(x_i)
    = \ln\!\left(\frac{x_i + \tfrac{1}{2}}{g(\mathbf{x} + \tfrac{1}{2})}\right),
  \label{eq:clr}
\end{equation}
where $g(\cdot)$ is the geometric mean across the taxon vector.
The half-count pseudocount ($+0.5$) was chosen over Laplace ($+1$) or
multiplicative alternatives because it is small enough to leave
proportions of well-represented taxa essentially undisturbed while
stabilising structural zeros; this convention is consistent with ALDEx2
and ANCOM-BC practice~\citep{gloor2017microbiome}.
In Python, Equation~\ref{eq:clr} reduces to:
\begin{center}\small\ttfamily
np.log(X + 0.5) - np.log(X + 0.5).mean(axis=1, keepdims=True)
\end{center}

\subsection{Global taxon-set selection}
\label{subsec:taxa_selection}

We selected the top 80 taxa by prevalence (the fraction of samples exceeding
each taxon's own median abundance) computed across \emph{all} 3,409 stool
samples before any group partitioning.
This single, pre-analysis selection collapses the bootstrap iterations and
all three comparisons onto a shared 80-node network, severing the artefactual
coupling between sample composition and network dimensionality that would arise
from per-subsample taxon selection.

\textbf{Prevalence definition.}
Standard presence--absence prevalence (the proportion of samples with a
detected read for taxon $j$) is not applicable after \clr\ transformation
because \clr\ maps structural zeros to large negative values rather than to
zero, making every taxon formally ``present'' in every sample.
We therefore define prevalence as the fraction of samples in which taxon $j$'s
\clr\ value exceeds its own column-wise median---equivalently, the 50th-percentile
exceedance rate.
This selects taxa that are consistently in the upper half of their within-taxon
abundance distribution, a natural analogue of detection-based prevalence that
avoids compositional distortion.
Sensitivity to the precise threshold (i.e., whether rank-selection at $N = 50$,
$80$, or $120$ changes the conclusions) is evaluated in
Section~\ref{subsec:sensitivity}.

\subsection{Network construction and Vietoris--Rips filtration}
\label{subsec:networks}

Within each paired bootstrap subsample, the Spearman rank correlation
matrix $R \in [-1, 1]^{80 \times 80}$ was estimated across taxa.
Entries were converted to a dissimilarity measure:
\begin{equation}
  d_{ij} = 1 - |r_{ij}|,
  \label{eq:dist}
\end{equation}
mapping strongly correlated taxa (positive or negative) to near-zero
distances and independent taxa to $d_{ij} = 1$.
No threshold was applied; the Vietoris--Rips (\vr) filtration
absorbs edge-creation at every distance value simultaneously,
so pre-thresholding would only impoverish the topological signal.

\subsection{Persistent homology computation}
\label{subsec:ph}

Persistent homology was computed with Ripser~\citep{tralie2018ripser}
(\texttt{ripser} v0.6.14) operating directly on the
$80 \times 80$ dissimilarity matrix.
We extracted homology groups $H_0$ (connected components) and $H_1$
(one-cycles, i.e., loops).
$H_2$ computation (two-dimensional voids, the filling of tetrahedra)
was omitted for two reasons.
First, three-dimensional co-occurrence motifs lack a clear biological
interpretation in microbial interaction networks; the scientifically
motivated hypothesis---that dysbiosis disrupts cross-feeding
\emph{cycles}---maps onto $H_1$, not $H_2$.
Second, constructing the 3-simplices required for $H_2$ increases both
computation and memory requirements substantially relative to the $H_1$
case; for the 80-node networks analysed here, $H_1$ alone already
saturates the topological signal at a tractable runtime.

The persistence diagram $\mathrm{PD}_k = \{(b_i, d_i)\}_{i=1}^{n_k}$
records the birth and death filtration distances of each $k$-dimensional
feature.
Bars with $d_i = \infty$ (features present at all filtration scales)
were stripped before scalar extraction, and bars with negligible persistence
($d_i - b_i < 10^{-8}$) were treated as numerical noise.

\subsection{Scalar topological features}
\label{subsec:features}

Six scalars were extracted from each $\mathrm{PD}_1$:

\begin{description}[leftmargin=2.2em, itemsep=3pt]
  \item[\textbf{H1 count.}] Cardinality $|\mathrm{PD}_1|$: raw number of
  detected co-occurrence cycles.
  \item[\textbf{H1 entropy.}] Persistent entropy~\citep{patania2017topological}:
  $\mathcal{S} = -\sum_i \frac{\ell_i}{L} \log \frac{\ell_i}{L}$,
  where $\ell_i = d_i - b_i$ and $L = \sum_i \ell_i$.
  Quantifies the evenness of cycle-lifetime distribution.
  \item[\textbf{Total persistence.}] $\sum_i \ell_i$: cumulative
  topological ``weight'' of all detected cycles.
  \item[\textbf{Mean lifetime.}] $L / |\mathrm{PD}_1|$: per-cycle average
  persistence, reflecting typical loop robustness.
  \item[\textbf{Max lifetime.}] $\max_i \ell_i$: the single most
  tenacious loop; robust to diagram sparsity.
  \item[\textbf{Max Betti-1.}] Peak value of the function
  $\beta_1(\varepsilon) = |\{i : b_i \leq \varepsilon < d_i\}|$:
  the maximum number of simultaneously active loops at any filtration
  threshold.
\end{description}

\subsection{Paired bootstrap design}
\label{subsec:bootstrap}

We ran $I = 200$ paired iterations for each comparison.
Iteration $t$ proceeds as follows:
\begin{enumerate}[leftmargin=*, itemsep=2pt]
  \item Sample $n_A = n_B = 100$ observations \emph{without} replacement
  from groups A and B respectively.
  \item Compute the six topological features independently for each
  subsample: $\mathbf{f}^{(t,A)}$ and $\mathbf{f}^{(t,B)}$.
  \item Record element-wise deltas $\boldsymbol{\delta}^{(t)}
    = \mathbf{f}^{(t,A)} - \mathbf{f}^{(t,B)}$.
\end{enumerate}
The observed test statistic for feature $k$ is $\bar{\delta}_k
= I^{-1} \sum_t \delta_k^{(t)}$.
Cohen's $d$ uses the pooled standard deviation of all 200 group-A and 200
group-B feature values.

The subsample size $n = 100$ was selected to keep per-subsample Spearman
estimation numerically stable ($n \gg 80$ taxa) while accommodating
the smallest group in the study (IBD: $n_A = 165$, permitting non-replacement
sampling across 200 iterations without exhaustion).

\subsection{Sign-flip label permutation}
\label{subsec:stats}

For each feature $k$, a null distribution is generated by independently and
uniformly randomizing the sign of each delta:
\begin{equation}
  T^{(\pi)} = \frac{1}{I}\sum_{t=1}^{I} \sigma_t^{(\pi)}\, \delta_k^{(t)},
  \qquad \sigma_t^{(\pi)} \overset{\text{i.i.d.}}{\sim} \{-1, +1\},
  \label{eq:signflip}
\end{equation}
repeated 500 times (yielding a minimum achievable $p = 1/501 \approx 0.002$,
sufficient to reach the floor at our observed significance levels).
Because subsamples overlap, iteration-level deltas are not strictly
independent; inference relies exclusively on sign-flip permutation, which
preserves this correlation structure under the null.
The permutation $p$-value is
$p = \bigl(|\{\pi : |T^{(\pi)}| \geq |\bar{\delta}_k|\}| + 1\bigr) / 501$.
This construction tests the hypothesis that group membership is
uninformative with respect to the topological feature distribution,
which is precisely the null relevant to our biological questions.
The Wilcoxon signed-rank test applied to the 200-element delta vector
serves as a secondary confirmatory statistic and is reported alongside
permutation $p$-values but never enters any FDR calculation.

\paragraph{FDR policy (pre-specified).}
Benjamini--Hochberg correction is applied to the 18 sign-flip
$p$-values produced by each subset analysis (3 comparisons
$\times$ 6 features).
Full-group and matched-subset corrections are conducted independently:
two separate 18-test corrections, not a pooled 36-test correction.
This is the \emph{only} multiple-testing adjustment applied; Wilcoxon
$p$-values are untouched.

\subsection{Covariate matching}
\label{subsec:matching}

Each comparison was replicated within a subsample stratified by three
demographic covariates: age (bins $<25$, 25--40, 40--60, $>60$\,yr), sex
(female/male; ambiguous responses excluded), and BMI ($<25$, 25--30,
$>30$\,kg\,m$^{-2}$).
Within each stratum, all group-A members were retained; group-B members
were sampled at up to $3 \times$ the stratum-A count.
Of the 3,409 stool samples, 2,740 (80.4\%) carried complete covariate records.
Matched group sizes are given in Table~\ref{tab:groups}.

\subsection{Taxa-size sensitivity}
\label{subsec:design_sensitivity}

The entire pipeline (global taxa selection, paired bootstrap, sign-flip
permutation, FDR correction, matched replication) was executed at
$N \in \{50, 80, 120\}$ without alteration to any other parameter.

\subsection{Classification benchmark}
\label{subsec:classification_methods}

To situate the topological features within the broader predictive-modelling
literature, we benchmarked nine feature sets against each other for
binary IBD classification (healthy vs.\ IBD-positive):
\begin{enumerate}[leftmargin=*, itemsep=2pt]
  \item Shannon diversity only.
  \item Topology only (6 $H_1$ scalars, per-sample k-NN neighbourhood TDA).
  \item Topology + Shannon.
  \item Aitchison PCoA(10): top-10 principal coordinates of the Euclidean
        distance on CLR-80.
  \item Bray--Curtis PCoA(10): top-10 principal coordinates of Bray--Curtis
        dissimilarity on relative abundances of the same 80 taxa.
  \item Aitchison PCoA + Topology.
  \item Bray--Curtis PCoA + Topology.
  \item Aitchison PCoA + Topology + Shannon.
  \item Bray--Curtis PCoA + Topology + Shannon.
\end{enumerate}

\paragraph{Nested PCoA (no test-fold leakage).}
For each cross-validation fold, PCoA was fitted on the training-fold distance
matrix only. Test samples were projected into the training eigenvector space
via Gower's~\cite{gower1968adding} out-of-sample formula, ensuring no
test-fold information enters the feature representation.

\paragraph{Cross-validation protocol.}
All classifiers were evaluated with 5-fold stratified cross-validation
(seed = 42).
Two classifier families were used: $L_2$-penalised logistic regression and
random forest (scikit-learn defaults).
Features were standardised within each fold using a \texttt{StandardScaler}
fitted on training samples only.
Performance is reported as area under the receiver-operating curve
(AUC-ROC), averaged across folds.

\paragraph{Per-sample topology and transductive design note.}
Per-sample $H_1$ features were computed using a $k$-nearest-neighbour
neighbourhood ($k = 60$) constructed from the full sample set prior to
cross-validation splitting.
Consequently, test-sample neighbourhoods include training samples.
This is a transductive rather than inductive design: labels are never
seen during feature computation, so there is no label leakage, but the
features are not strictly out-of-sample.
All AUC values for topology-containing feature sets should be interpreted
under this constraint.

\subsection{Implementation and reproducibility}
\label{subsec:repro}

The pipeline runs on Python\,3.11 with \texttt{ripser}\,v0.6.14,
\texttt{persim}, \texttt{scipy}\,1.17, \texttt{numpy}\,2.4, and
\texttt{pandas}.
The random seed was fixed to 42 throughout.
Source code, download scripts, result CSVs, and figure generation are
archived at
\href{https://github.com/adam-jfkhs/Microbiome-TDA}%
{\texttt{github.com/adam-jfkhs/Microbiome-TDA}};
the relevant entry points are \texttt{scripts/run\_agp\_bootstrap\_v2.py}
and \texttt{scripts/run\_taxa\_sensitivity.py}.
